\chapter{Conclusion}
This thesis presents \textit{glaciVis}, an interactive visualization tool designed to support data preselection for machine learning applications in climate and glacier research.
By combining structured preprocessing with coordinated spatial and temporal views, the tool enables informed feature assessment prior to model development.

During the development process, the tool proved to be well suited for gaining deep insights into meteorological data.
It allows users to precisely assess whether a dataset is suitable for a given task and, if so, which temporal resolution and spatial extent are necessary. The use case in Chapter~\ref{ch:use_case} illustrated this concretely: the exploration of wind variables revealed that the available 10\,m components in ERA5-Land are unlikely to carry predictive information for glacier mass balance, while
suggesting that other ERA5 products may contain more relevant variables.

The project demonstrates that interactive visualization can serve as a critical intermediate step between raw data acquisition and automated learning methods.
It provides transparency in feature selection and supports hypothesis driven exploration.